\chapter{Review of Related Works}
\label{ch:rrl}

\section{Digital Transformation in Higher Education Administration}
\label{sec:digital-transformation}

The transition from traditional administrative processes to digital ecosystems is a critical focus for modern educational institutions. Early Management Information Systems (MIS) successfully digitized basic student records but frequently functioned as siloed databases that failed to integrate cross-departmental workflows \cite{ref3,ref5}. This fragmentation resulted in underutilized resources, compromised administrative efficiency, and unequal educational access, not merely because paper was replaced by screens, but because the underlying workflows remained disconnected. Institutions globally have recognized that streamlining student services requires unified digital workflows as a structural prerequisite for operational sustainability \cite{ref46}. Contemporary digital transformation therefore requires interoperable architectures: systems capable of unifying data collection, processing, and reporting across departments into a single, coherent workflow \cite{ref36}. Kayanja et al. \cite{ref46} confirm in their 2025 study of fully automated and paperless systems in higher education that the most significant efficiency gains emerge not from digitizing isolated records, but from architecturally integrating those records into cross-departmental process flows. This finding is directly applicable to scholarship administration, where document submission, verification, interview scoring, and committee deliberation currently exist as separate, disconnected processes at most Philippine higher education institutions. At AdDU, these steps span the University Scholarship Office, interview panels, five School Scholarship Subcommittees, and the Office of Admission. The literature therefore establishes that the effectiveness of digital transformation in higher education depends on interoperable, cross-departmental architectures, the foundational design principle motivating ScholarPath AdDU's centralized, Supabase-backed infrastructure.

\section{Web-Based Scholarship Management Systems}
\label{sec:web-systems}

While web technologies are increasingly leveraged for scholarship management globally, the literature focusing on Southeast Asian and Philippine higher education contexts reveals both what existing systems achieve and where they fall short. Bagaporo and Espinosa \cite{ref12} developed an Online Scholarship Management System for Camarines Sur Polytechnic Colleges (CSPC), implementing a centralized digital repository that consolidates scholarship listings, automates announcement dissemination to enrolled students, and provides an administrative dashboard for monitoring scholarship status. The system successfully eliminated the need for physical bulletin boards and improved information reach across the multi-campus institution. However, the platform stops at information dissemination: document submission remained a manual, offline process, and students were still required to process physical paperwork and submit it to the respective offices in person \cite{ref12}.

Daluyon and Bilog \cite{ref14} developed a Scholarship Management Information System specifically for State Universities and Colleges (SUCs) in Rizal, incorporating online application submission and a structured administrative review interface. Their system advanced beyond the CSPC implementation by digitizing the submission process, reducing paper routing, and enabling administrators to view and sort applications from a centralized dashboard. Crucially, however, Bilog's evaluation found that unaided self-assessment of eligibility was the primary driver of application errors and administrative rework across multi-campus deployments. Students frequently submitted applications for grants they did not qualify for, generating review bottlenecks for administrative staff \cite{ref14}. This finding motivates ScholarPath AdDU's Phase~1 Pre-qualification, in which the platform flags baseline criteria early and a University Scholarship Office staff member confirms the outcome, rather than leaving applicants to judge their own eligibility or letting software decide it.

Al-Ayyubi and Maulana \cite{ref3} designed a web-based scholarship management information system for a higher education institution in Indonesia, achieving efficient dissemination of scholarship information and basic application form routing through a unified web interface. The system demonstrated measurable improvements in processing speed and information accessibility. Yet the researchers themselves identified the architectural absence of integrated, secure document handling as a critical and recurring gap: applicants entered details through forms, but prerequisite documents such as transcripts and income proofs were still submitted physically and managed separately \cite{ref3}. This gap informs the Document Vault architecture in ScholarPath AdDU, which receives documents in two submission phases and supports correction within the platform.

At the international level, Amer \cite{ref5} at the Lebanese American University developed a Smart Scholarship Management System as an all-in-one solution, integrating application submission, document handling, and status tracking into a unified platform. The system demonstrated the technical feasibility of end-to-end digital scholarship administration but was built for a single decision-making body applying one set of criteria. The challenge of serving several academic units that each apply their own selection rules to a shared applicant pool was not a design constraint of that system. The reviewed literature therefore reveals a clear progression in capability: from static digital repositories \cite{ref12}, to submission-enabled portals \cite{ref14,ref3}, to integrated multi-function platforms \cite{ref5}. None demonstrates a validated solution that combines phased digital document submission, structured correction of incomplete files, and decision support for multiple school-level committees whose every decision is recorded, the combination that ScholarPath AdDU addresses.

\section{Automated Decision-Making versus Decision Support in Scholarship Selection}
\label{sec:decision-support}

In digital scholarship systems, the technical mechanism by which applicants are evaluated against eligibility criteria varies substantially across the reviewed literature, and these differences in mechanism determine both the quality of outcomes and who is accountable for them. Three dominant implementation patterns emerge: keyword-overlap matching, threshold-based sequential rule evaluation, and machine-learning-driven classification, each carrying distinct computational behaviors and institutional tradeoffs. Keyword-overlap matching operates by tokenizing profile fields and tallying term coincidences against grant descriptions. Systems relying solely on this mechanism are computationally incapable of evaluating intersecting numerical constraints as a unified logical gate. Orgianus et al. \cite{ref36} document this explicitly: keyword-based retrieval produced persistent false-positive matches because a keyword engine performs term frequency comparisons, not constraint satisfaction. Threshold-based sequential rule evaluation, as implemented by Bilog \cite{ref14}, advances beyond keyword matching by encoding eligibility criteria as conditional logic, but evaluates constraints independently and without formal priority ordering, producing conflict resolution failures at edge cases that Section~\ref{sec:screening-deployed} examines in detail. A third pattern, machine-learning classification, is documented by Kanwal et al. \cite{ref50}, whose C4.5 decision tree achieved 95.62\% prediction accuracy on scholarship eligibility data. However, ML-derived decision boundaries carry two institutional liabilities: the model requires retraining whenever criteria change, and its decisions offer no interpretable audit trail, a concern for accountability and for compliance under RA 10173 that the authors themselves acknowledge \cite{ref50}.

These findings establish a hierarchy of tradeoffs: keyword matching fails at constraint evaluation; sequential rule evaluation fails at conflict resolution; ML classification fails at interpretability and update responsiveness. A common thread runs through all three. Each pattern places the eligibility or selection decision inside the software, and each failure mode therefore becomes a wrong decision made without a responsible person. ScholarPath AdDU draws the opposite conclusion from the same evidence. Its rules encode baseline and school-specific criteria as screening flags and sort orders, but the decision that changes an applicant's status is always made and recorded by a University Scholarship Office staff member, an interview panel, or a School Scholarship Subcommittee. This design is consistent with the institutional directive, arising from AdDU's preparation for a quality audit, that no algorithm selects scholars \cite{ref52}. % TODO(L-6): add team-supplied sources on decision support systems / human oversight of automated decisions (ref53+).

\section{Search and Filtering Mechanisms for Committee Review}
\label{sec:filtering-rrl}

To understand why faceted filtering is a suitable retrieval architecture for committee review of a large applicant pool, it is necessary to examine how keyword search operates at a structural level. Keyword search systems maintain, at their technical core, an inverted index: a data structure that maps individual terms to the ordered list of records in which they appear (posting lists). When a user submits a query, the engine retrieves the posting lists for each query term and intersects or unions them, ranking results by term frequency or relevance weighting \cite{ref43}. This mechanism performs correctly when query vocabulary aligns with document vocabulary, but a keyword engine is structurally incapable of filtering by continuous numerical attributes. It cannot natively express ``show applicants in this school's partition with a verified high school average of at least 85\% and a complete Phase~2 file,'' because inverted indexes store term occurrences, not numerical bounds or states. Mahdi et al. \cite{ref4} confirm this limitation in their review of 47 deployed information retrieval systems, finding that keyword-only retrieval consistently underperforms in repositories where items carry multiple intersecting categorical and numerical attributes. Faceted filtering resolves this structurally: the system exposes its underlying metadata as interactive filter controls, allowing users to iteratively reduce a large result set by selecting attribute values \cite{ref43}. The query mechanism shifts from posting-list intersection over terms to predicate conjunction over structured metadata fields.

The performance characteristics of this approach are tightly coupled to the indexing architecture underlying the query executor. For categorical attributes with low cardinality, such as school, track, or lifecycle status, B-Tree indexes support O(log n) key lookups and range scans that retrieve qualifying rows without full-table sequential scans. For continuous numerical and temporal attributes, such as high school average, interview score, and the start of a Conditional Revert window, Range indexes support boundary-constrained retrieval \cite{ref36}. When compound indexes are constructed across combinations of these fields, the query planner can satisfy multi-facet filter combinations from a single traversal rather than executing independent scans per facet and joining their results in memory. Tunkelang \cite{ref43} establishes that the quality of the underlying metadata taxonomy is the decisive determinant of whether this architecture reduces or merely redistributes cognitive load: poorly normalized facets produce overlapping filter categories that reintroduce the ambiguity the search was designed to eliminate. The reviewed literature does not document this architecture applied to committee review of a scholarship applicant pool that is partitioned across several academic units, each with its own criteria. ScholarPath AdDU applies B-Tree and Range compound indexes over school, track, program choice, status, verification, and score facets, with client-side input debouncing at 300 milliseconds.

\subsection{Policy Desynchronization in Manual Systems}
\label{sec:desync}

A critical vulnerability of decentralized, paper-based administrative systems is policy desynchronization, the structural mismatch between the rules as published and the rules as applied \cite{ref16}, compounded by the multi-campus application error patterns documented by Bilog \cite{ref14}. When criteria or deadlines change, static notices and separately maintained files have no programmatic mechanism to propagate those changes to dependent records. The data model underlying these systems is the root cause: criteria and deadlines are stored as unstructured content rather than as structured, queryable records in a relational database. A change to a threshold is therefore a content editing operation, not a data update, and carries no guarantee of consistency across all instances where that criterion is referenced. Catibog \cite{ref16} identifies this structural weakness as a primary driver of discrepancies between published criteria and the criteria actually applied at award time. At AdDU, the risk is multiplied by five School Scholarship Subcommittees that each apply their own criteria to a shared applicant pool. A centralized, database-driven platform in which the SOP and each school's rule set are stored once and referenced everywhere is architecturally positioned to resolve this, but validated implementations of this model remain absent from the reviewed Philippine HEI literature.

\section{Eligibility Screening and Committee Workflows in Deployed Systems}
\label{sec:screening-deployed}

While Section~\ref{sec:decision-support} examined where eligibility decisions should reside and Section~\ref{sec:filtering-rrl} examined filtering for committee review, it is necessary to examine how these technologies have been implemented in deployed systems. In the Philippine context, Bilog \cite{ref14} implemented a rule-based validation layer within the Scholarship Management Information System for SUCs in Rizal, where eligibility constraints, including GPA thresholds, institutional affiliation, and enrollment status, were encoded as conditional logic applied during application intake. While this implementation reduced erroneous submissions at the point of entry, it evaluated constraints sequentially and independently, without a mechanism for handling conflicts when multiple rules applied simultaneously. This limitation resulted in edge cases where students holding concurrent government grants passed initial validation for additional state-sponsored awards \cite{ref14}. ScholarPath AdDU responds to this conflict-handling gap differently from a rule engine that resolves conflicts itself: when rules conflict, the platform raises every applicable flag and presents them together to the responsible committee, which records the decision.

At the international level, Amer \cite{ref5} implemented an eligibility filtering component within the Smart Scholarship Management System at the Lebanese American University that matched student academic profiles against scholarship criteria defined in a centralized database. The system used a multi-criteria filter that evaluated grade point average, program of study, and nationality. However, the filter operated on a flat, non-hierarchical rule set, and a single body applied it to the whole applicant pool \cite{ref5}. ScholarPath AdDU's screening instead operates across five school-specific rule sets, such as the SON zero-conditional policy and the SEA entrance exam cutoffs, applied to the partitions of a shared applicant pool.

Regarding filtering, Mahdi et al. \cite{ref4} identify compound indexing and real-time metadata synchronization as the dominant performance optimization strategies in large-scale faceted deployments, and Tunkelang \cite{ref43} establishes that faceted search effectiveness depends on a well-normalized taxonomy. These findings inform the Committee Portal's normalized facets and compound indexes. Taken together, the reviewed implementations establish that rule-based screening and faceted filtering are validated approaches in educational portal contexts, but no prior system has combined school-partitioned rule-based screening with human-recorded committee decisions in a single platform. This implementation gap is what ScholarPath AdDU is positioned to address.

\section{Accessibility and Remote Access in Web-Based Systems}
\label{sec:accessibility}

Remote access in web-based academic systems is not merely a deployment preference; it is an architectural decision that determines whether a system can sustain adoption in mobile-dominant, bandwidth-variable environments such as Mindanao's higher education landscape. The technical architecture a system adopts for cross-platform delivery determines its responsiveness under variable connectivity, its ability to deliver proactive notifications, and its maintenance overhead, all of which have direct consequences for adoption rates and operational sustainability. Three dominant architectural patterns exist for cross-platform academic portals, each representing a distinct tradeoff between development complexity, performance characteristics, and native capability access. Responsive web applications (RWAs) adapt a single web codebase to varying viewport sizes through CSS media queries and fluid grid layouts, requiring no separate build for mobile targets but limited to capabilities available through the browser's JavaScript API surface. Progressive Web Applications (PWAs) extend this pattern by registering Service Workers, background scripts that cache application shells and API responses in the browser's Cache Storage, enabling partial offline functionality and reducing cold-start load times under intermittent connectivity \cite{ref33}. However, iOS imposes persistent platform restrictions on PWA background execution and push notification delivery through WebKit that reduce their reliability as proactive alert channels; applications with time-critical deadline and correction notices cannot depend on this model for consistent delivery. Hybrid mobile architectures, specifically the Capacitor framework used in ScholarPath AdDU, resolve this by compiling the same Vite production JavaScript bundle into native Android and iOS packages: Capacitor instantiates a WKWebView (iOS) or WebView (Android) host process and bridges a set of native plugin APIs, including push notifications via APNs and FCM, local filesystem access, and camera, into the JavaScript execution context through a typed bridge interface. This architecture delivers the development efficiency of a unified JavaScript codebase while accessing native OS capabilities that browser-based delivery cannot reach, without requiring a parallel Swift or Kotlin codebase.

The functional case for this architectural choice is reinforced by documented adoption evidence specific to the Philippine context. The World Bank's \cite{ref48} assessment of digital transformation in Philippine higher education explicitly identifies mobile-first design as a structural prerequisite for adoption, not a secondary usability consideration, particularly for Mindanao-based institutions where smartphone access substantially exceeds laptop ownership rates. For scholarship applicants, who photograph and upload documents such as report cards and certificates, a phone camera connected directly to the Document Vault is the most direct submission channel. Yoo and Huang \cite{ref49} further establish that platforms perceived as culturally or technologically misaligned with users' existing digital habits generate adoption failure through reversion: when a portal's interface model conflicts with how a user's primary device renders and interacts with web content, users abandon the digital channel and return to manual processes, recreating the administrative bottlenecks the platform was deployed to eliminate. What prior scholarship portal implementations in the Philippine literature have not documented is a formal operationalization of the Capacitor-wrapped production build approach, producing a single deployable codebase for web, Android, and iOS, within a scholarship admission context. ScholarPath AdDU's hybrid architecture addresses this gap, enabling the notification subsystem to deliver time-critical deadline and Conditional Revert notices through native push channels.

\section{Usability and User Experience in Educational Platforms}
\label{sec:usability-rrl}

In evaluating web-based academic systems, usability testing is universally identified as the most vital assessment method. Task-based usability testing directly measures how well users can complete real tasks under realistic conditions. To complement observed performance, researchers rely on standardized perception-based instruments like the System Usability Scale (SUS). Bangor, Kortum, and Miller \cite{ref13} conducted an extensive empirical evaluation of the SUS, confirming its high reliability and validity across diverse digital interfaces. Empirical research by Tullis and Stetson \cite{ref42} further demonstrated that a sample size of 10 to 12 participants is sufficient to yield reliable SUS scores that accurately reflect the broader user experience. When these standardized instruments are used alongside the ISO/IEC 25010 Software Quality Evaluation model \cite{ref31}, developers can effectively quantify perceived usability and ensure the interface minimizes administrative friction. The evaluative methodology of this study is built on this established consensus: task-based usability testing paired with the SUS and the ISO/IEC 25010 quality framework constitutes the field's most validated approach for determining whether an academic platform genuinely reduces the friction it was designed to address \cite{ref13,ref31,ref42}. Applying this protocol to ScholarPath AdDU provides the empirical backbone through which the study's claims about submission effort, lifecycle correctness, committee usefulness, and deadline management are substantiated.

\section{Policy, Data Privacy, and Access Control in Financial Aid Information Systems}
\label{sec:privacy-rrl}

Scholarship admission systems process data that is simultaneously sensitive, multi-tenanted, and multi-role: an applicant's household income record must be readable by University Scholarship Office staff but invisible to other applicants; an interview panel must see only the applications assigned to it; a School Scholarship Subcommittee must see only its own school's partition. The technical mechanism by which these constraints are enforced, not merely the policy that states they should be enforced, determines whether a system achieves genuine RA 10173 compliance or only nominal compliance. The two dominant access control paradigms in information systems differ fundamentally in their enforcement architecture. Access Control Lists (ACLs) enforce access by maintaining, for each data object, an explicit list of subjects and their permitted operations. In a scholarship context, this means maintaining per-user permission entries for each application record, a model that scales as O(users $\times$ objects) and requires manual permission reassignment each time a user's scope is modified. More consequentially, ACL enforcement in web applications is typically implemented at the application layer: a middleware check verifies the requesting user's permission before executing the database query. This model carries a structural security vulnerability: if an API endpoint is inadvertently exposed through a misconfigured route, a CORS policy error, a parameter injection attack, or a client-side logic bypass, the database itself executes the query without resistance \cite{ref38}.

Role-Based Access Control (RBAC), formalized by Ferraiolo and Kuhn \cite{ref24} and standardized through NIST \cite{ref35}, resolves the scalability problem through role indirection: permissions are attached to roles rather than to individual users, reducing permission management from O(users $\times$ objects) to O(roles $\times$ permissions). The more architecturally significant property of RBAC for RA 10173 compliance, however, is its capacity for enforcement at the database engine level through PostgreSQL's Row Level Security (RLS) mechanism. An RLS policy is a predicate expression, for example \texttt{USING (school\_code = auth\_school())} on the Applications table for subcommittee members, that PostgreSQL evaluates as part of every query's execution plan before any rows are returned. Because this predicate is evaluated within the query executor itself, it operates below the application layer entirely. Despite this technically superior enforcement model, Halili \cite{ref26} documents that Philippine HEI scholarship portals commonly lack transparent, auditable access control implementations, and CHED's memorandum guidelines \cite{ref18} do not specify the technical mechanism through which access segregation must be achieved. The consequence is that existing Philippine scholarship portals may be formally RA 10173-compliant at the policy documentation level while remaining technically vulnerable at the implementation level. ScholarPath AdDU addresses this gap by enforcing its RBAC permission matrix through PostgreSQL RLS policies across all system roles: applicants (self-access only), University Scholarship Office staff (all applications, for verification), interview panels (assigned applications only), School Scholarship Subcommittees (own school's partition only), the Office of Admission (accepted lists), and the System Administrator (configuration only).

\section{Theoretical Framework}
\label{sec:framework}

This study is anchored on the Information Systems Success Model \cite{ref20}, a widely validated framework that evaluates the effectiveness of centralized web portals. The model posits that the success of an information system is driven by six interrelated dimensions: System Quality, Information Quality, Service Quality, System Use, User Satisfaction, and Net Benefits. As illustrated in Figure~\ref{fig:issm}, this standard model has been adapted to map the operational variables and institutional constraints of the ScholarPath AdDU platform.

\begin{figure}[htbp]
  \centering
  % TODO(FIG): redraw figure1_issm.png -- the current image still shows
  % "Automated Matching", "Accurate QPI data", "Prompt OSA support", and
  % "Equitable Aid Distribution". See context/iso_revision_tracker.md.
  \includegraphics[width=0.85\textwidth]{figure1_issm.png}
  \caption{Adapted Information Systems Success Model for ScholarPath AdDU}
  \label{fig:issm}
\end{figure}

In the context of this study, the foundational dimensions are defined by the platform's technical and informational reliability. System Quality is represented by the portal's 24/7 accessibility, its secure data architecture, and an audit log that traces every status change to a person, each of which responds to the system-architecture gaps documented in the reviewed literature \cite{ref3,ref12,ref14}. Information Quality is reflected in the accuracy of verified applicant data, school-partitioned committee views, and real-time status notifications, addressing the policy desynchronization and information fragmentation identified in \cite{ref16} and \cite{ref36}. Service Quality encompasses the system's responsiveness, including clear correction notices for applicants and prompt workflow support for the University Scholarship Office, informed by the cultural alignment requirements documented in \cite{ref49} and \cite{ref48}. According to the model, high performance in these foundational areas drives user engagement. As applicants, staff, and subcommittee members experience reliable system and information quality, their System Use (e.g., uploading documents, verifying files, recording decisions) and User Satisfaction (measured quantitatively via the System Usability Scale) will increase. Ultimately, this sustained engagement leads to the system's Net Benefits: reduced administrative bottlenecks, fewer applications lost to missing documents, and committee-led decisions that are traceable in an institutional audit. This entire framework operates within the parameters of institutional policy and national data privacy law, with all dimensions of system success governed by the compliance requirements of RA 10173 through the RBAC architecture established in \cite{ref24,ref38}.

\section{Synthesis of RRL and Research Gaps}
\label{sec:synthesis}

The reviewed literature, taken as a whole, establishes a clear progression of scholarship management system capability, from information dissemination \cite{ref12} to submission-enabled portals \cite{ref3,ref14} to multi-function integrated platforms \cite{ref5}, while consistently revealing that no existing implementation resolves the full set of challenges posed by a multi-committee scholarship admission process. The specific gaps motivating ScholarPath AdDU are summarized in Table~\ref{tab:synthesis}, with each gap traced to the reviewed work that documented it.

\begin{table}[htbp]
  \centering
  \small
  \caption{Synthesis of System Gaps and Proposed Design Responses}
  \label{tab:synthesis}
  \begin{tabular}{@{}p{3.0cm}p{5.4cm}p{6.0cm}@{}}
    \toprule
    Reviewed Work & System/Technology Gap Identified & ScholarPath AdDU Design Response \\
    \midrule
    Al-Ayyubi \& Maulana \cite{ref3}; Bagaporo \& Espinosa \cite{ref12} & Portals stop at information dissemination; document submission remains physical and disconnected from the digital workflow. & Two-stage digital submission into the Document Vault, with phase-tagged metadata and versioned corrections under Conditional Revert. \\
    Bilog \cite{ref14} & Unaided eligibility self-assessment drives application errors; rule-based filtering lacked conflict handling. & Phase~1 Pre-qualification flags baseline criteria early; conflicting flags are shown together to staff or the subcommittee, who record the decision. \\
    Kanwal et al. \cite{ref50} & Machine-learned eligibility decisions offer no interpretable audit trail. & Data-Feeding Committee Portal: no module changes an applicant's status without a recorded human actor; every change is written to an audit log. \\
    Amer \cite{ref5} & Single decision body with one flat rule set; multi-unit selection not addressed. & Applicant pool partitioned by school, with five School Scholarship Subcommittee rule sets presented as flags and sort orders. \\
    Mahdi et al. \cite{ref4}; Tunkelang \cite{ref43} & Keyword search fails over multi-attribute records; facets require normalized taxonomy and compound indexing. & Committee filtering with B-Tree and Range compound indexes over school, track, program, status, verification, and score facets; 300\,ms debounce. \\
    Catibog \cite{ref16}; Orgianus et al. \cite{ref36} & Policy desynchronization between published and applied criteria in decentralized systems. & Centralized Supabase/PostgreSQL backend storing the SOP and each school's rule set once; updates propagate to all partitions. \\
    Yoo \& Huang \cite{ref49}; World Bank \cite{ref48} & Misaligned interfaces drive users back to manual processes; mobile-first design is a prerequisite for adoption. & Hybrid web and mobile architecture via Capacitor, from a single Vite codebase. \\
    Halili \cite{ref26}; CHED \cite{ref18} & Limited documentation of RBAC operationalization under RA 10173 in multi-role workflows. & RBAC matrix enforced through PostgreSQL RLS for applicants, staff, interview panels, subcommittees, the Office of Admission, and the System Administrator. \\
    \bottomrule
  \end{tabular}
\end{table}

The convergence of these gaps, namely disconnected document workflows, unaided or automated eligibility decisions without accountability, single-body selection models, policy desynchronization, misaligned interfaces, and undocumented RBAC compliance in Philippine HEI contexts, identifies a distinct and unaddressed design space. No existing implementation in the reviewed literature simultaneously integrates phased digital document submission with structured correction, school-partitioned decision support for multiple selection committees, an audit trail that traces every status change to a person, and proactive notification within a single platform. This gap directly motivates the present study and grounds each architectural decision in ScholarPath AdDU in specific, documented evidence from the reviewed literature.
